\subsection{可一致化空间}

在第二章，§ 4，第 1 号中，我们提出了刻画可一致化拓扑空间的问题。其解答由下述定理给出：

定理 2. 拓扑空间 X 可一致化的必要充分条件是它满足下述公理：

\((\mathrm{O}_{\mathrm{IV}})\) 对任意点 \(x_0 \in X\) 与任意邻域 \(V\) （ \(x_0\) 的），存在 \(X\) 上的一个连续实值函数，它取值于 \([0,1]\) ，在 \(x_0\) 处等于 0 ，且在 CV 上等于 1 。

该条件是必要的。因为若 X 上存在一个与 X 的拓扑相容的一致结构，则由定理 1，这个一致结构可以由 X 上的一族伪度量 \((f_{t})\) 定义，并且我们可以不失一般性地假定这个族是饱和的（第 2 号）。由这样的伪度量族所定义的一致结构的围域的定义，存在一个伪度量 \(f_{\alpha}\) （族 \((f_{t})\) 中的）以及一个数 \(a > 0\) ，使得 \(f_{\alpha}(x_0, x) \geqslant a\) 对一切 \(x \in \mathbb{C}\mathrm{V}\) 成立。由此推出，函数 \(g(x) = \inf \left(1, \frac{1}{a} f_{\alpha}(x_0, x)\right)\) 满足 \((\mathrm{O}_{\mathrm{IV}})\) 中规定的所有条件。

该条件是充分的。因为设 \(\Phi\) 是 X 到 {[}o, i{]} 中的所有连续映射组成的集。公理 (O \(_{IV}\) ) 表明，使属于 \(\Phi\) 的所有函数都一致连续的最粗的一致结构与 X 的拓扑相容（第二章，§ 2，第 3 号）。

定义 4. 拓扑空间称为完全正则的，如果它可一致化且是豪斯多夫的。

由定理 2，等价地，一个空间是完全正则的，如果它满足公理 (H) {[}参看第一章，§ 8，第 1 号，命题 1{]} 与 \((\mathbf{O}_{\mathbf{IV}})\) 。

注。公理 \((\mathrm{O}_{\mathrm{IV}})\) 蕴含 \((\mathrm{O}_{\mathrm{III}})\) （参看第一章，§ 8，第 4 号），因为若 V 是 \(x_0\) 的一个邻域，且 \(f\) 是 X 上取值于 {[}o, r{]} 的连续实值函数，满足 \(f(x_0) = o\) ，\(f(x) = 1\) 对一切 \(x \in \complement V\) ，则集 \(\overline{f}([0, 1/2])\) 是包含在 V 中的 \(x_0\) 的一个闭邻域。特别地，每个完全正则空间都是正则的（这说明了该术语的合理性）。但也存在不是完全正则的正则空间的例子 (*)，因而 \((\mathrm{O}_{\mathrm{III}})\) 不蕴含 \((\mathrm{O}_{\mathrm{IV}})\) 。

每个紧空间都是完全正则的（第二章，§ 4，第 1 号，定理 1），因此紧空间的每个子空间也都是完全正则的。我们现在可以通过证明其逆命题来完善这一命题：

命题 3. 拓扑空间 X 完全正则的必要充分条件是它同胚于某个紧空间的一个子空间。

考察 X 上使所有从 X 到 {[}o, 1{]} 中的连续映射都一致连续的最粗的一致结构；在定理 2 的证明中我们已经用过这个一致结构，并且看到当 X 可一致化时它与 X 的拓扑相容。此外，由区间 {[}o, 1{]} 的紧性以及第二章，§ 4，第 2 号的命题 3，这个一致结构是一个预紧空间的结构。于是若 X 是豪斯多夫的，则 X 关于这个一致结构的完备化是紧的，命题得证。

因此我们可以说，完全正则空间可以嵌入到一个紧空间中。把这一结果表述成下面的形式常常是方便的：

一般地，方体是指拓扑空间 \(K^{I}\) ，即一族拓扑空间的乘积，族中每个拓扑空间都等同于 R 的一个紧区间 K，其指标集为任意集合 I。若 I 有限且有 n 个元素，我们就重新得到在第六章，§ 1，第 1 号中定义的 n 维闭立方体的概念。方体是紧空间（第一章，§ 9，第 5 号，定理 3）。

命题 4. 若拓扑空间 X 完全正则，则它同胚于某个方体的一个子空间。

用 \((f_{i})_{i\in \mathbf{I}}\) 表示 \(\mathbf{X}\) 到 \(\mathbf{K} = [\mathrm{o},\mathrm{i}]\) 中的所有连续映射组成的族，并用 \(g\) 表示映射 \(x\to (f_i(x))\) （\(\mathbf{X}\) 到

\(K^{I}\) 中的）。若 x, y 是 X 中任意两个不同的点，则由公理 (H) 与 \((\mathrm{O}_{\mathrm{IV}})\) 可知存在一个指标 t，使得 \(f_{t}(x) \neq f_{t}(y)\) ，因而 g 是 X 到 \(K^{I}\) 中的一个单射。此外，直接可知：g 是使所有 \(f_{t}\) 都一致连续的 X 上最粗的一致结构到 \(g(\mathbf{X})\) 上由 \(K^{I}\) 的乘积一致结构诱导的一致结构上的一个同构；更有甚者，g 是 X 到 \(g(\mathbf{X})\) 上的一个同胚。

